%% ==========================================================================
%%  Approach 6 --- a pure goods reformulation.  OPEN, and the current best lead.
%%
%%  Source: derived 2026-08-04.  Scripts: update_6/targetG.py, targetG_adv.py.
%%
%%  STATUS.  Target G implies Conjecture 1 (proved).  Target G itself is open,
%%  and survives the adversarial generator that refuted Approach 5 in one run,
%%  including on the very instance that did the refuting.
%%
%%  WHY THIS SECTION EXISTS.  Corollary cor:coverage says coverage IS the whole
%%  gap, so the replica framing is an equivalence and not a reduction --
%%  inventing there means inventing on the original problem with extra
%%  bookkeeping.  This section moves instead to the coordinate system in which
%%  the literature's algorithms already live.
%% ==========================================================================
\section{Approach 6: A Pure Goods Reformulation}
\label{sec:approach6}

Every approach so far has worked inside the chore model or the replica model.
This one leaves both. The observation is that the size-shift transform of
Theorem~\ref{thm:sizeshift} is not merely a comparison device: it is an
involution of the dichotomous class onto itself, so the chore problem can be
restated as a question about \emph{goods} with no chores, no replica, no
coverage constraint and no schedule anywhere in it.

% --------------------------------------------------------------------------
\subsection{The transform is an involution}

\begin{lemma}
\label{lem:sizeshift-bijection}
The map $\cost \mapsto \tilde{v}$, $\tilde{v}(S) := \abs{S} - \cost(S)$, is an
involution of the set of dichotomous set functions (\Cref{def:dichcost}) onto
itself.
\end{lemma}

\begin{proof}
$\tilde{v}(\emptyset) = 0$, and for $g \notin S$ the marginal is
$\tilde{v}(S \cup \set{g}) - \tilde{v}(S) = 1 - \big(\cost(S \cup \set{g}) -
\cost(S)\big) \in \set{0,1}$, so $\tilde{v}$ is dichotomous. Applying the map
twice returns $\abs{S} - (\abs{S} - \cost(S)) = \cost(S)$.
\end{proof}

So every dichotomous \emph{goods} instance arises as $\tilde{v}$ for exactly one
dichotomous \emph{chore} instance and conversely: the two models are the same
object in two coordinate systems, and nothing is lost by working in either.

% --------------------------------------------------------------------------
\subsection{The target}

Fix an allocation $\alloc$ and write $\tilde{\subsidy}$ for the minimum subsidy
vector of the goods instance $\set{\tilde{v}_i}$ at $\alloc$, which is defined
exactly when the chore instance is envy-freeable at $\alloc$, since by
Theorem~\ref{thm:sizeshift} the two envy graphs have the same cycle weights.
Put
\[
  q_i \;:=\; \tilde{\subsidy}_i + \abs{\bundle{i}}
  \qquad\text{---\; goods subsidy plus bundle size.}
\]

\begin{conjecture}[Target G --- \textbf{open}]
\label{conj:targetG}
Every dichotomous goods instance admits an allocation whose vector $q$ has
spread at most $1$, i.e.\ $\max_i q_i - \min_i q_i \le 1$.
\end{conjecture}

\begin{proposition}
\label{prop:targetG-implies}
Conjecture~\ref{conj:targetG} implies Conjecture~\ref{conj:main}.
\end{proposition}

\begin{proof}
By Theorem~\ref{thm:sizeshift},
$\pathw{\alloc}(u) = \max_v [\tilde{d}_{\alloc}(u,v) + \abs{\bundle{u}} -
\abs{\bundle{v}}]$, where $\tilde{d}_{\alloc}(u,v)$ is the heaviest $u \to v$
path in the goods envy graph. Appending to any $u \to v$ path a heaviest path
leaving $v$ gives $\tilde{d}_{\alloc}(u,v) + \tilde{\subsidy}_v \le
\tilde{\subsidy}_u$, so $\tilde{d}_{\alloc}(u,v) \le \tilde{\subsidy}_u -
\tilde{\subsidy}_v$ and therefore
\[
  \pathw{\alloc}(u) \;\le\; \max_v\big[(\tilde{\subsidy}_u + \abs{\bundle{u}})
  - (\tilde{\subsidy}_v + \abs{\bundle{v}})\big] \;=\; q_u - \min_v q_v .
\]
A spread of at most $1$ therefore gives $\pathw{\alloc} \le 1$, and
Observation~\ref{obs:integrality} upgrades this to
$\optsubsidy \in \set{0,1}^n$ with total at most $n-1$.
\end{proof}

% --------------------------------------------------------------------------
\subsection{Why this is a better place to work}

Three things distinguish Target G from everything preceding it.

\textbf{It asks existing algorithms for one more property, rather than asking
for a new algorithm.} The two ingredients of $q$ are exactly the two guarantees
the literature already delivers separately. \cite{BKNS22} produces
$\tilde{\subsidy} \in \set{0,1}^n$ on dichotomous goods; \cite{BDNSV20}
and~\cite{LMS26} produce \emph{balanced} allocations, the latter by padding with
dummies so that every agent takes exactly one object per round, giving
$\abs{\abs{\bundle{i}} - \abs{\bundle{j}}} \le 1$. Neither controls the other
quantity. Target G asks for both at once, with the residual requirement being a
\emph{correlation}: among balanced allocations, the agents holding the larger
bundles must be the ones needing no subsidy. That is intuitively the right way
round --- in the goods instance a larger bundle means more value --- and it is a
concrete ask of machinery that already exists.

\textbf{It has no scheduling and no coverage.} Approach~3 fails on reachability
through a state space with dead ends and Approach~4 on rules fixed before the
allocation is known. Target G quantifies over allocations directly, so neither
failure mode is available to it.

\textbf{It is arc-free in the sense Remark~\ref{rem:arc-vs-path} demands.} The
quantity $\tilde{\subsidy}_i$ is itself a longest-path weight, so Target G is a
condition on paths, not on arcs. That is exactly what killed Approach~5, and
Target G is stated in terms the obstruction cannot reach.

% --------------------------------------------------------------------------
\subsection{Evidence, tested against the adversary that killed Approach~5}

Uniform balance also passed exhaustive $n = m = 3$ and thousands of random
instances before dying on the first structured adversarial run, so Target G was
put through the same adversary rather than the uniform sampler.

\begin{center}
\begin{tabular}{@{}p{0.52\textwidth}p{0.36\textwidth}@{}}
\toprule
Test & Target G failures \\
\midrule
Exhaustive, $n = m = 3$, all $9{,}880$ instances & $0$ \\
The discrepancy counterexample of
Proposition~\ref{prop:no-balance}, which refuted Approach~5 & $0$ \\
The insertion obstruction witness, the no-go instance of
Theorem~\ref{thm:w4}, and the instance of
Proposition~\ref{prop:msw-false} & $0$ \\
Exhaustive over structured families, $m \le 4$, $n = 3$
($20{,}436$ instances) & $0$ \\
Endpoint-constant sweeps to $n = 5$, $m = 5$ ($4{,}520$ instances) & $0$ \\
Randomised, $n \le 5$, $m \le 5$ ($3{,}400$ instances) & $0$ \\
\bottomrule
\end{tabular}
\end{center}

At $n = m = 3$ the best achievable spread is $0$ on $9{,}001$ instances and $1$
on the remaining $879$ --- never $2$. That the majority admit spread exactly
zero, meaning $\tilde{\subsidy}_i + \abs{\bundle{i}}$ can be made \emph{constant}
across agents, is a stronger regularity than Target G needs and is the first
thing a proof attempt should try to explain.

\begin{remark}[What would falsify it, and what it would not settle]
\label{rem:targetG-caveats}
Target G is sufficient for Conjecture~\ref{conj:main} but not known to be
necessary, since the inequality $\tilde{d}_{\alloc}(u,v) \le \tilde{\subsidy}_u
- \tilde{\subsidy}_v$ used in Proposition~\ref{prop:targetG-implies} can be
slack. So a counterexample to Target G would kill this route without touching
the conjecture --- the same relationship Remark~\ref{rem:converse} records for
the peel frame, and the reason both must be tested separately from the
conjecture itself.

The proof obligation is correspondingly precise. Take any algorithm producing a
balanced allocation on the goods instance --- iterated maximum-weight perfect
matching is the obvious candidate --- and show that its output can be chosen so
that no agent holding a maximum-size bundle requires a subsidy. Equivalently:
show that a heaviest path in the goods envy graph never starts at an agent with
a largest bundle. That is a statement about one algorithm on one model, with no
chores in it.
\end{remark}
